
%===================================================================================================%


\frame{\frametitle{Testing strategies}
Here are some ideas to check consistency of code output.
I think of them as Cross Checks.
Some are the \textbf{same checks you do in physics}.
What I use most commonly (no particular order, not exhaustive, not disjoint):
\begin{enumerate}
\item Compare against trusted result
% can be from literature, another code base, etc 
\item Compare with an analytic calculation
\item Compute a result multiple ways
\item Check for consistency between observables
% for example if one obs is the derivative of the other,
% that constrains how they must be related
\item Check known properties
% probabilities in [0,1]; 2d ising magnetization monotonic in T
\item Verify independence of hyperparameters
% dependence on jackknife bin size? solve tolerance? 
\end{enumerate}
}

\frame{\frametitle{CC1: Compare against trusted result}
\begin{itemize}
\item Often a starting point: \textbf{reproduce something}
% e.g. if you are implementing something new following the method
% of a paper, start by reproducing what's in that paper
\item Danger: you \textbf{assume other calculation is correct}
\begin{itemize}
\item Has the result already been reproduced?
\item You could apply some Cross Checks to the trusted result
\item If you are unsure about trusted result,
\textbf{blinding} can help against human bias: 
\begin{enumerate}
\item Let $r\in[-R,R]$ unknown to you 
\item Rescale your observable by $r$
\item Remove $r$ when you are confident in your calculation
\end{enumerate}
\end{itemize}
\item Can test against your old results to \textbf{ensure stability}
\begin{itemize}
\item Some of your unit tests could work this way
\end{itemize}
\end{itemize}
% Example of CMD? Or maybe the electron whatever thing
}

\frame{\frametitle{CC2: Compare with an analytic calculation}
\begin{itemize}
\item Simple test case \textbf{calculable by hand}
\item Known result from theory
\end{itemize}
% diff_deriv
}


\frame{\frametitle{CC3: Compute a result multiple ways}
There's usually more than one way to solve a problem
\begin{itemize}
\item Check optimized implementation against easy-to-understand one
\item Test different algorithms against each other
\item Test different models against each other
\end{itemize}
% Example of correlator class in GPU
% Different integration techniques give same result
}


\frame{\frametitle{CC4: Check for consistency between observables}
Two observables may be constrained by relationships between them.
Design tests that ensure observables obey these known constraints.
% Acceleration/velocity?
% What kinds of checks did Frithjof do?
% Maybe one of your checks from material params
}


\frame{\frametitle{CC5: Check known properties}
Depending on an observable's interpretation, you may already
know some properties it must have. Test them! Examples:
\begin{enumerate}
\item is your covariance matrix positive semidefinite?
\item is your unitary matrix actually unitary?
\item is your RHMC reversible?
\item is your probability between 0 and 1?
\end{enumerate}
}


\frame{\frametitle{CC6: Verify independence of hyperparameters}
\begin{itemize}
\item Parameter: Directly \textbf{model-relevant} number
\begin{itemize}
\item Temperature, volume, particle mass, coupling constant
\item Fit parameters 
\end{itemize}
\item Hyperparameter: Needed \textbf{only for computational strategy}
\begin{itemize}
\item Jackknife bins, bootstrap samples
\item Initial guess of a fitter
\item Max iterations, solve tolerance, CNN kernel size 
\end{itemize}
\item Often only independent of hyperparameter \textbf{in some region}
\begin{itemize}
\item Number of jackknife bins until error bar plateaus 
\item Initial guesses leading to reasonable $\chi^2/{\rm d.o.f.}$
\end{itemize}
\end{itemize}
}



